\subsection{Total elimination depth}
\label{sec:total-depth}

The same layered graph $S_q(G)$ gives a gap for total elimination
depth, also called VPT-sum~\cite{HogemoEtAl2021}. Here every child
contributes to the objective, and each vertex of the common prefix
increases the depth of every vertex below it. We reuse the
first-branch structure and bound these contributions together.

For a rooted forest $F$, let $\iota(F)$ be the number of unordered
pairs of distinct vertices that are incomparable in its ancestor order.
This counts the savings in total depth relative to one chain.

\begin{lemma}[Incomparable-pair accounting]
\label{lem:incomparable-pair-accounting}
For every rooted forest $F$ on $N$ vertices, the first identity below
holds. If $F$ is nonempty and not a single chain, define its prefix
$P$ and children $F_1,\ldots,F_\nu$ as in
Section~\ref{subsec:ugc-branch-notation}, and put $N_j=|V(F_j)|$
for $j\in[\nu]$. Then the second identity holds as well:
\[
 \totd(F)=\frac{N(N+1)}2-\iota(F),\qquad \iota(F)=\sum_{1\leq i<j\leq\nu}N_iN_j+
           \sum_{j=1}^{\nu}\iota(F_j).
\]
Adding a chain above all roots leaves $\iota(F)$ unchanged.
\end{lemma}

\begin{proof}
A vertex's depth counts itself and its proper ancestors. Thus total
depth counts all $N$ vertices and all comparable unordered pairs,
proving the first identity, including $N=0$. Every prefix vertex
is comparable with every vertex, while vertices in different children
are incomparable. This partitions the incomparable pairs as stated.
New vertices on a chain above all roots are comparable with every
vertex, so they create no incomparable pair.
\end{proof}

\subsubsection{The YES upper bound}

The elimination forest from Lemma~\ref{lem:ugc-completeness} has a
particularly simple depth sum. Its common prefix contains all copies
of the vertices outside an independent set. The remaining copies form
one chain per layer.

\begin{lemma}[Total depth of the cover-first forest]
\label{lem:total-depth-completeness}
For every graph $G$ on $n$ vertices and every integer $q\geq1$,
\[
 \totd(S_q(G))\leq
 \frac{qn(qn+1)-q(q-1)\alphaI(G)^2}{2}.
\]
\end{lemma}

\begin{proof}
Choose a maximum independent set of $G$, and use the forest $F$ of
Lemma~\ref{lem:ugc-completeness}. Its only incomparable pairs have endpoints in
different lower chains. There are $q$ such chains, each with
$\alphaI(G)$ vertices, so there are
$\binom q2\alphaI(G)^2$ incomparable pairs. Therefore the forest has
total depth
\[
 \totd(F)\overset{\text{Lem.~\ref{lem:incomparable-pair-accounting}}}{=}
 \frac{qn(qn+1)}2-\binom q2\alphaI(G)^2
 =\frac{qn(qn+1)-q(q-1)\alphaI(G)^2}{2},
\]
which gives the claimed upper bound, also when the common prefix is empty.
\end{proof}

\subsubsection{A structural lower bound for total depth}

Use the notation from Section~\ref{subsec:ugc-branch-notation}.
Lemma~\ref{lem:incomparable-pair-accounting} accounts for the prefix
through one pair-counting identity: only incomparable pairs save depth. Restoring a child's
missing pairs in $U_j\times I_j$ above its roots preserves this
saving and lets induction apply to a full layered graph.

The next inequality combines the squared child terms and their
pairwise products. Cauchy--Schwarz will then bound its right-hand side
using the total layer count and the sum of the child vertex-cover numbers.

\begin{lemma}[Quadratic branching inequality]
\label{lem:quadratic-branching}
For a positive integer $\nu$ and nonnegative reals $W_1,\ldots,W_\nu$,
\[
 \sum_{j=1}^{\nu} W_j^2+14\sum_{1\leq i<j\leq \nu}W_iW_j
 \leq\left(\sum_{j=1}^{\nu} \sqrt{W_j}\right)^4.
\]
\end{lemma}

Its proof, by expansion into nonnegative terms, is given in
Appendix~\ref{app:total-depth-branching}.

\begin{theorem}[Structural total-depth lower bound]
\label{thm:total-depth-structural}
For every finite simple graph $G$ on $n$ vertices and every integer
$q\geq1$,
\begin{equation}\label{eq:total-depth-structural}
 \totd(S_q(G))\geq
 \frac37q^2\vc(G)^2+\frac{qn}{2}.
\end{equation}
\end{theorem}

\begin{proof}
We induct on $q$, with the induction hypothesis asserted for every
finite simple source graph. If $n=0$, both sides are zero. Fix
$n\geq1$ and an arbitrary elimination forest $F$ of $S_q(G)$.
A single-chain forest has total depth
$qn(qn+1)/2\ge(3/7)q^2\vc(G)^2+qn/2$, since $\vc(G)\le n$ and
$1/2\ge3/7$. This includes $q=1$, since $S_1(G)$ is a clique.
Otherwise use the first-branch notation
from Section~\ref{subsec:ugc-branch-notation}. Put $c=\vc(G)$ and,
for $j\in[\nu]$, put $c_j=\vc(G[U_j])$, $W_j=q_jc_j$, and $L_j=q_jn_j$.
Also put $S=\sum_{j=1}^{\nu} W_j$, $A=q\alphaI(G)$, and $N=qn$.
We have $n\overset{\eqref{eq:cover-independent-complement}}{=}c+\alphaI(G)$
and therefore $N=qc+A$. The shared child resource bounds give
\begin{equation}\label{eq:total-depth-child-resources}
 \left(\sum_{j=1}^{\nu}q_j\right)
 \left(\sum_{j=1}^{\nu}c_j\right)
 \overset{\eqref{eq:ugc-child-source-bounds}}{\leq}qc.
\end{equation}
Consequently,
\begin{equation}\label{eq:total-depth-support-bounds}
S=\sum_{j=1}^{\nu} q_jc_j
   \leq\Bigl(\sum_{j=1}^{\nu} q_j\Bigr)\Bigl(\sum_{j=1}^{\nu} c_j\Bigr)
   \overset{\eqref{eq:total-depth-child-resources}}{\leq}qc,\qquad \sum_{j=1}^{\nu} L_j=\sum_{j=1}^{\nu} q_jn_j
   \overset{\eqref{eq:ugc-child-source-bounds}}{\leq}
   S+\alphaI(G)\sum_{j=1}^{\nu} q_j
   \overset{\eqref{eq:ugc-child-source-bounds}}{\leq}S+A.
\end{equation}

For each $j\in[\nu]$, restore the $M_j=L_j-N_j$ missing vertices of child $j$ above
its roots, obtaining a forest $\widetilde F_j$ of
$S_{q_j}(G[U_j])$ by Lemma~\ref{lem:restore-total-depth}.
Since $q_j\overset{\text{Lem.~\ref{lem:ugc-layer-usage}}}{<}q$,
induction bounds its total depth below by $3W_j^2/7+L_j/2$.
Restoration preserves incomparable pairs, so
\begin{equation}\label{eq:total-depth-child-bound}
 \iota(F_j)
 \overset{\text{Lem.~\ref{lem:incomparable-pair-accounting}}}{=}
 \frac{L_j(L_j+1)}2-\totd(\widetilde F_j)
 \overset{\text{induction}}{\leq}\frac{L_j^2}{2}-\frac37W_j^2.
\end{equation}
The additional incomparable pairs lie in different children.
Since $N_j\leq L_j$, their contribution and the squared terms
in~\eqref{eq:total-depth-child-bound} satisfy
\begin{equation}\label{eq:total-depth-cross-child-loss}
\sum_{1\leq i<j\leq\nu}N_iN_j+\frac12\sum_{j=1}^{\nu} L_j^2
 \overset{\text{Lem.~\ref{lem:ugc-child-restriction}}}{\leq}
 \sum_{1\leq i<j\leq\nu}L_iL_j+\frac12\sum_{j=1}^{\nu} L_j^2 =\frac12\Bigl(\sum_{j=1}^{\nu} L_j\Bigr)^2
 \overset{\eqref{eq:total-depth-support-bounds}}{\leq}\frac{(S+A)^2}{2}.
\end{equation}
Subtracting all incomparable pairs from the chain cost yields
\begin{equation}\label{eq:total-depth-combined-children}
\begin{aligned}
 \totd(F)&\overset{\text{Lem.~\ref{lem:incomparable-pair-accounting};
 \eqref{eq:total-depth-child-bound}}}{\geq}
 \frac{N^2+N}{2}-\sum_{1\leq i<j\leq\nu}N_iN_j
             -\frac12\sum_{j=1}^{\nu} L_j^2+\frac37\sum_{j=1}^{\nu} W_j^2\\
 &\overset{\eqref{eq:total-depth-cross-child-loss}}{\geq}
 \frac{N^2+N}{2}-\frac{(S+A)^2}{2}+\frac37\sum_{j=1}^{\nu} W_j^2.
\end{aligned}
\end{equation}
The terms involving only the $W_j$ can be grouped as
\[
 \frac{S^2}{2}-\frac37\sum_{j=1}^{\nu}W_j^2
 =\frac1{14}\left(\sum_{j=1}^{\nu}W_j^2
                  +14\sum_{1\leq i<j\leq\nu}W_iW_j\right).
\]
Apply Lemma~\ref{lem:quadratic-branching}, then Cauchy--Schwarz using
$\sqrt{W_j}=\sqrt{q_j}\sqrt{c_j}$:
\begin{equation}\label{eq:total-depth-quadratic-estimate}
\frac{S^2}{2}-\frac37\sum_{j=1}^{\nu} W_j^2
 \overset{\text{Lem.~\ref{lem:quadratic-branching}}}{\leq}
 \frac1{14}\left(\sum_{j=1}^{\nu} \sqrt{W_j}\right)^4 \overset{\text{Cauchy--Schwarz}}{\leq}
 \frac1{14}\left(\sum_{j=1}^{\nu} q_j\right)^2\left(\sum_{j=1}^{\nu} c_j\right)^2
 \overset{\eqref{eq:total-depth-child-resources}}{\leq}\frac{q^2c^2}{14}.
\end{equation}
Since $N=qc+A$ and $A\geq0$, the bound $S\leq qc$ gives
\[
 N^2-(S+A)^2=q^2c^2-S^2+2A(qc-S)
 \overset{\eqref{eq:total-depth-support-bounds}}{\geq}q^2c^2-S^2.
\]
Substituting this into~\eqref{eq:total-depth-combined-children}
leaves only the quadratic estimate:
\[
 \totd(F)\overset{\eqref{eq:total-depth-combined-children},\,
 \eqref{eq:total-depth-support-bounds}}{\geq}
 \frac N2+\frac{q^2c^2}{2}
 -\left(\frac{S^2}{2}-\frac37\sum_{j=1}^{\nu} W_j^2\right)
 \overset{\eqref{eq:total-depth-quadratic-estimate}}{\geq}
 \frac N2+\frac{q^2c^2}{2}-\frac{q^2c^2}{14}
 =\frac37q^2c^2+\frac N2.
\]
\end{proof}

\subsubsection{Sharpness of the structural coefficient}

The coefficient $3/7$ cannot be increased in a bound of the form
\eqref{eq:total-depth-structural} valid for every $G$ and $q$.
The following forests approach equality even when the source
independence number is a vanishing fraction of its order.

\begin{proposition}\label{prop:total-depth-structural-sharpness}
Let $k\geq0$ and $t\geq2$ be integers, put $q=2^k$, and let $G$
be the disjoint union of $q$ copies of $K_t$. Write $n=qt$ and
$N=qn$. There is an elimination forest $F$ of $S_q(G)$ with
\[
 \totd(F)=\frac N2+
 \left(\frac37+\frac{1}{14\cdot8^k}\right)N^2.
\]
Consequently, if a constant $\beta$ satisfies
$\totd(S_q(G))\geq\beta q^2\vc(G)^2+qn/2$ for every finite simple
graph $G$ and every integer $q\geq1$, then $\beta\leq3/7$.
\end{proposition}

% \begin{proof}
% We construct the forest recursively. For $k=0$, the output is
% $K_t$, and a chain has total depth $t(t+1)/2$.
% For $k\geq1$, partition the layer indices into equal sets $I_1,I_2$
% and the source clique-components into equal groups, whose vertex
% sets are $U_1,U_2$. Put all vertices in
% $(U_1\times I_2)\cup(U_2\times I_1)$ on a common prefix.
% Its length is $N/2$. Below this prefix put separate forests on
% $U_1\times I_1$ and $U_2\times I_2$, each of order $N/4$.
% No output edge joins these two sets: their layers differ and no
% source edge joins $U_1$ to $U_2$. Each induced child graph has
% $q/2$ layers over a disjoint union of $q/2$ copies of $K_t$, so
% the recursive construction applies. Every edge incident with the
% prefix has comparable endpoints, proving that the forest is valid.
%
% For fixed $t$, let $T_k$ be its total depth and $N_k=4^kt$ its
% order. The prefix contributes its own depths and $N_k/2$ to the
% depth of each of the $N_k/2$ child vertices. Thus
% \[
% \begin{aligned}
%  T_k
%  &\overset{\text{Lem.~\ref{lem:ugc-total-depth-accounting}}}{=}
%  \frac{(N_k/2)(N_k/2+1)}2+\frac{N_k^2}{4}+2T_{k-1}\\
%  &=\frac{3N_k^2}{8}+\frac{N_k}{4}+2T_{k-1}.
% \end{aligned}
% \]
% Writing $T_k=N_k/2+a_kN_k^2$ and using $N_{k-1}=N_k/4$ gives
% \[
%  a_0=\frac12,\qquad
%  a_k=\frac38+\frac{a_{k-1}}8,\qquad
%  a_k-\frac37=\frac18\left(a_{k-1}-\frac37\right).
% \]
% Iterating the last identity yields
% $a_k=3/7+1/(14\cdot8^k)$, proving the formula.
%
% Here $\alphaI(G)=q$ and $\vc(G)=q(t-1)$. The proposed universal
% bound and the constructed forest would imply
% \[
%  \beta
%  \overset{\text{assumed bound}}{\leq}
%  \frac{\totd(S_q(G))-N/2}{q^2\vc(G)^2}
%  \overset{\totd(S_q(G))\leq\totd(F)}{\leq}
%  \left(\frac37+\frac{1}{14\cdot8^k}\right)
%  \frac{t^2}{(t-1)^2}.
% \]
% Letting $t$ tend to infinity and then $k$ tend to infinity gives
% $\beta\leq3/7$. Also $\alphaI(G)/n=1/t$, as claimed.
% \end{proof}

This sharpness statement concerns the universal structural bound;
it does not rule out stronger hardness from additional source
properties or a different construction.

\subsubsection{The approximation gap}
\label{subsec:ugc-total-depth-gap}

\begin{proof}[Dense total-depth gap: proof sketch]
Lemma~\ref{lem:total-depth-completeness} and
Theorem~\ref{thm:total-depth-structural} give zero-error quadratic
coefficients $q(3q+1)/8$ in the YES case and $3q^2/7$ in the NO
case, with the same linear term $qn/2$. Their ratio
$24q/(7(3q+1))$ tends to $8/7$. For a fixed $1\leq r<8/7$,
choose a fixed $q$ and then a small
source error so that the quadratic coefficients have ratio greater
than $r$. A fixed number of disjoint source copies
makes the quadratic gap dominate the linear term while preserving
the independent-set promise and maximum degree.
Appendix~\ref{app:dense-ugc-total} gives the coefficient expansion,
parameter choices, and explicit replication count.
\end{proof}
